\chapter{Maintenance and updates}
\label{chap:mantenimiento}

\section{Save a known working configuration}

After the first successful run on the board, save the basic environment information.
It provides a reference if an update causes a fault.

\begin{lstlisting}[language=bash]
mkdir -p evidence/known-good

akd1000-bringup doctor --json \
  > evidence/known-good/doctor.json

python -m pip freeze \
  > evidence/known-good/pip-freeze.txt

sha256sum models/*.fbz \
  > evidence/known-good/model-hashes.txt
\end{lstlisting}

Also retain the project and driver versions used. If the system later stops working,
this configuration helps identify what has changed.

\section{Update the environment}

Updates to the kernel, driver, Python or Akida can affect board operation. Where
possible, change one component at a time and repeat the smoke test afterwards.

A simple procedure is:

\begin{enumerate}
  \item check compatibility with the new version;
  \item apply the update;
  \item reboot if needed;
  \item repeat board detection and the software and hardware tests.
\end{enumerate}

If a fault occurs, compare the system with the saved configuration and follow the
procedure in \cref{chap:diagnostico}.

The driver uses DKMS to rebuild the module after a kernel update. This works only
while the new kernel version remains supported by the driver.

\section{Shutdown and storage}

Before disconnecting the board, finish active runs and shut down the host normally.
Then disconnect the power supplies before removing the board or FFC cable.

For long-term storage, protect the boards against electrostatic discharge and avoid
bending the FFC cable or leaving it under mechanical tension.
